\documentclass[11pt,a4paper]{article}
\usepackage[utf8]{inputenc}
\usepackage[margin=0.75in]{geometry}
\usepackage{amsmath,amssymb,amsfonts}
\usepackage{booktabs}
\usepackage{graphicx}
\usepackage{hyperref}
\usepackage{caption}
\usepackage{microtype}
\usepackage{xcolor}
\usepackage{placeins}
\usepackage{pdflscape}
\usepackage{float}
\usepackage{makecell}
\usepackage{longtable}

\hypersetup{
    colorlinks=true,
    linkcolor=blue!70!black,
    citecolor=blue!70!black,
    urlcolor=blue!70!black
}

\captionsetup[figure]{position=top, labelfont=bf, font=small, skip=4pt}
\captionsetup[table]{position=top, labelfont=bf, font=small, skip=4pt}

\title{\textbf{DISSERTATION PROGRESS REPORT} \\[0.5em] \Large \textbf{Chapter 2: From the Okishio Viability Condition to a Testable Capacity Equation (Specification B5)}}
\author{\textbf{Diego Polanco} \\ Department of Economics, UMass Amherst \\ \textit{Progress Report for Michael Ash and Dissertation Committee}}
\date{July 2026}

\begin{document}

\maketitle

\begin{abstract}
\noindent Marxian theories of technical change claim that the productive potential of a capital stock is not an engineering constant but the outcome of a cost-minimizing choice made under distributive and price pressure. If that claim holds, the terms through which distribution enters a capacity equation are structural parameters rather than controls. This report derives a cointegrating capacity equation from the Okishio viability condition formalized in Section 3.2 of this chapter, estimates it by Integrated Modified OLS over 1929--2024 under the E10V estimator lock, and reports the full centering triad rather than a single preferred model. \textbf{The central result is negative: the theory imposes two sign restrictions and no estimated specification satisfies both at once.} The uninteracted model recovers the wage-savings dividend ($\hat{\gamma}_1 = 4.20 > 0$) and reverses the relative price penalty ($\hat{\gamma}_2 = -1.00 < 0$); the break-interacted model recovers the price penalty ($\hat{\gamma}_2 = 0.56 > 0$) and reverses the wage dividend ($\hat{\gamma}_1 = -9.69 < 0$). The extensive scale elasticity is sharply identified at $\hat{\theta}_{NRC} = 0.588$ ($t = 8.09$), varying 1.8 percent across centering regimes, so the failure is specific to the state-dependent channel. A wage-mechanization structural break at 1969, reported as significant by an earlier specification, does not survive a correction to the interaction span ($t = 1.31$).
\end{abstract}

\section{Introduction}

Marxian theories of technical change make a sharp claim about capacity: the productive potential of a capital stock is not an engineering constant but the outcome of a cost-minimizing choice made under distributive and price pressure. If that claim is right, the terms through which distribution enters a capacity equation are structural parameters, not controls. Most empirical work treats them as controls, and the resulting interaction terms then look like specification search.

This paper examines whether the Okishio viability condition, formalized in Section 3.2 of this chapter, licenses a specific cointegrating capacity equation for the United States over 1929--2024, and whether that equation survives estimation. I derive the estimating equation from the firm's cost-surplus problem rather than positing it, estimate it by Integrated Modified OLS under the E10V estimator lock, and report the full centering triad rather than a single preferred model. The central result is negative and worth stating immediately: \textbf{the theory imposes two sign restrictions, and no estimated specification satisfies both at once.} The uninteracted model recovers the wage-savings dividend ($\hat{\gamma}_1 = 4.20 > 0$) but reverses the relative price penalty ($\hat{\gamma}_2 = -1.00 < 0$); the break-interacted model recovers the price penalty ($\hat{\gamma}_2 = 0.56 > 0$) but reverses the wage dividend ($\hat{\gamma}_1 = -9.69 < 0$). The mechanism the theory describes is not jointly identified in this data.

Two subsidiary findings support that reading. The extensive scale parameter $\theta_{NRC}$ is sharply and stably estimated at 0.59 ($t = 8.09$), varying by 1.8 percent across centering regimes, so the failure is specific to the state-dependent channel and not general to the design. And the wage-mechanization structural break at 1969, which an earlier specification reported as significant, does not survive a correction to the interaction span ($t = 1.31$).

Notation used throughout: $y^{NFC}_t$ is log real non-financial corporate value added; $k^{NRC}_t$ and $k^{ME}_t$ are log real stocks of nonresidential structures and of machinery and equipment; $\kappa_t \equiv k^{ME}_t - k^{NRC}_t$ is the \textbf{mechanization gap}; $p^{\kappa}_t \equiv \ln P^{ME}_t - \ln P^{NRC}_t$ is the log \textbf{relative price} of machinery to structures; and $\omega_{t-1}$ is the lagged wage share.

\FloatBarrier
\section{Descriptive Grounding}

The raw series carry a fact that constrains any account of the mechanization decision. Table~\ref{tab:descriptive} reports mean annual changes; for the log levels these are growth rates, and for $\kappa_t$ and $p^{\kappa}_t$, which are log ratios, they are divergence rates between two series.

\begin{table}[H]
\centering
\caption{Mean Annual Change, by Era}
\label{tab:descriptive}
\small
\begin{tabular}{p{7.0cm}p{2.6cm}p{2.6cm}p{2.6cm}}
\toprule
Series & Full sample 1929--2024 & Fordist 1929--1973 & Post-Fordist 1974--2024 \\
\midrule
Corporate value added ($y^{NFC}$), \% p.a. & 3.89 & 4.93 & 3.12 \\
Nonresidential structures ($k^{NRC}$), \% p.a. & 5.76 & 5.46 & 5.72 \\
Machinery and equipment ($k^{ME}$), \% p.a. & 6.21 & 6.95 & 5.30 \\
Mechanization gap ($\kappa$), pp p.a. & 0.45 & \textbf{1.49} & \textbf{-0.43} \\
Relative price ($p^{\kappa}$), pp p.a. & -0.69 & \textbf{0.90} & \textbf{-2.03} \\
Wage share ($\omega$), pp p.a. & -0.10 & -0.05 & -0.17 \\
\bottomrule
\multicolumn{4}{p{15.5cm}}{\footnotesize Note: Growth rates are mean annual log differences. For $\kappa$ and $p^{\kappa}$, which are log ratios, the entry is the annual divergence rate between the two component series, in percentage points. $\omega$ is a share, so its entry is the annual change in percentage points of the share. Source: \texttt{specB\_descriptive\_growth\_rates.csv}.}
\end{tabular}
\end{table}

Read the last two rows together. During Fordism the relative price of machinery \textbf{rose} 0.90 points a year while the mechanization gap \textbf{widened} 1.49 points a year. After 1974 the relative price \textbf{collapsed} at 2.03 points a year while the gap \textbf{narrowed} 0.43 points a year. Machinery became dramatically cheaper relative to structures, and firms responded by slowing the accumulation of machinery relative to plant.

That co-movement is positive in both eras, and Section~\ref{sec:theory} predicts it should be negative. The puzzle is visible in the data before any estimator touches it, which is why it belongs here rather than in a robustness appendix. Three figures under \texttt{output/figures/} show the same relationship non-parametrically: the three-panel structural decomposition, the mechanization drift series, and the relative price drift series.

Output growth decelerated by 1.8 points between eras while structures accumulation held nearly constant at about 5.7 percent. The extensive margin kept expanding as its capacity yield fell.

\FloatBarrier
\section{Theory: The Okishio Viability Condition}
\label{sec:theory}

Section 3.2 of this chapter derives capacity from a decentralized cost-minimization problem. Potential output decomposes into an extensive envelope and an intensive yield generated by the gap:

\begin{equation}
y^p_t = \theta_{NRC} k^{NRC}_t + \Psi_t(\kappa_t, \omega_{t-1}, p^{\kappa}_t)
\end{equation}

The firm chooses $\kappa_t$ against a mechanization frontier $\Phi(\kappa_t) = \lambda_0 + \lambda_1 \kappa_t + \lambda_2 \kappa_t^2$, maximizing the net viability surplus

\begin{equation}
\max_{\kappa_t} \ c(\kappa_t) = \Phi(\kappa_t) - \kappa_t \cdot MC(\omega_{t-1}, p^{\kappa}_t)
\end{equation}

where $MC$ is the marginal viability cost of widening the gap. Two channels set it. Raising $\kappa_t$ substitutes machinery for labor, and the value of that displacement is proportional to the lagged wage share, so a higher wage share \textbf{reduces} the net cost of mechanizing. Raising $\kappa_t$ also requires accumulating machinery faster than structures, at a cost proportional to the relative price, so a higher $p^{\kappa}_t$ \textbf{increases} it:

\begin{equation}
MC(\omega_{t-1}, p^{\kappa}_t) = \gamma_0 - \gamma_1 \omega_{t-1} + \gamma_2 p^{\kappa}_t, \qquad \gamma_1 > 0, \ \gamma_2 > 0
\end{equation}

Here $\gamma_1$ is the wage-savings dividend and $\gamma_2$ the relative price penalty. The first-order condition gives the optimal gap

\begin{equation}
\kappa_t^{*}(\omega_{t-1}, p^{\kappa}_t) = \frac{(\gamma_0 - \lambda_1) - \gamma_1 \omega_{t-1} + \gamma_2 p^{\kappa}_t}{2\lambda_2}
\end{equation}

with strict concavity $\lambda_2 < 0$ guaranteeing an interior maximum. The comparative statics are the induced-innovation mechanism:

\begin{equation}
\frac{\partial \kappa_t^{*}}{\partial \omega_{t-1}} = -\frac{\gamma_1}{2\lambda_2} > 0, \qquad \frac{\partial \kappa_t^{*}}{\partial p^{\kappa}_t} = \frac{\gamma_2}{2\lambda_2} < 0
\end{equation}

A rising wage share forces cost-minimizing firms to widen the gap. A rising relative price of machinery suppresses it.

\paragraph{What this means at the workplace.} The algebra above compresses a concrete institutional process. Capital becomes capacity through a conversion that happens on the shop floor: a given stock of structures and machines yields potential output only at some rate of throughput, set by shift patterns, crewing levels, maintenance regimes and the pace of work. The mechanization gap is the margin on which firms adjust that conversion. Widening it --- adding machines faster than floor space --- raises output per worker-hour and lets the firm economize on a wage bill that has become expensive. Distribution enters the capacity equation because the wage share is what makes that substitution worth undertaking, and it enters \textit{lagged} because the wage settlement precedes the investment decision it provokes. The relative price enters because the same substitution has to be bought at the going price of equipment. Neither variable is a control for something else; each is a term in the calculation a cost-minimizing firm actually performs.

\paragraph{An unresolved conflict in the theory.} Section 3.2 line 68 derives $\partial\kappa^{*}/\partial p^{\kappa} < 0$. The companion note \texttt{A06\_RelativePrice\_Mechanization\_Theory.md} line 72 derives the opposite, $\partial q^{*}/\partial p^{rel} > 0$, from a multiplicative rather than additive cost form. The two are not reconciled in the vault. The descriptive evidence above sits with A06. This report records the conflict and does not adjudicate it; the ledger \texttt{specB5\_price\_channel\_sign\_conflict.csv} carries both derivations against every estimate.

\FloatBarrier
\section{How the Theory Licenses the Specification}

The step from theory to an estimating equation is an envelope-theorem argument, and it is what distinguishes this specification from a regression with interaction terms added for flexibility.

The marginal capacity payoff of the gap is the derivative of the intensive yield, which by the envelope theorem equals the net marginal viability yield at the optimum:

\begin{equation}
MP_{\kappa} = \frac{\partial \Phi}{\partial \kappa_t} - MC(\omega_{t-1}, p^{\kappa}_t) = \underbrace{(\lambda_1 - \gamma_0 + 2\lambda_2 \kappa_t)}_{\text{autonomous yield}} + \gamma_1 \omega_{t-1} - \gamma_2 p^{\kappa}_t
\end{equation}

Integrating over $\kappa_t$, holding the slow-moving state variables predetermined, gives the intensive yield and therefore the capacity equation:

\begin{equation}
y^p_t = \theta_{NRC} k^{NRC}_t + \theta_{\kappa} \kappa_t + \theta_{\omega\kappa}(\omega_{t-1}\kappa_t) + \theta_{\kappa p\kappa}(p^{\kappa}_t \kappa_t)
\label{eq:b5}
\end{equation}

with the structural identities

\begin{equation}
\theta_{\omega\kappa} = \gamma_1 > 0, \qquad \theta_{\kappa p\kappa} = -\gamma_2 < 0
\end{equation}

The interaction terms are the macroeconomic footprint of the viability constraint. Distribution and relative prices enter multiplicatively because they alter the marginal capacity payoff of the gap, not because a level effect needed conditioning.

Four consequences bind the estimation.

\textbf{The standalone gap is required.} $\theta_{\kappa}$ carries the autonomous yield $(\lambda_1 - \gamma_0)$. An earlier specification (B4) dropped it on the argument that mechanization is ``only a channel.'' The derivation does not support that: without $\theta_{\kappa}$ the marginal payoff has no intercept and the transformation elasticity cannot be recovered from the structural model. B5 restores it.

\textbf{There is no three-way term.} $MC$ is additively separable in $\omega_{t-1}$ and $p^{\kappa}_t$, so the model generates two two-way interactions and nothing else. A $\kappa\omega p^{\kappa}$ term would invent structure the theory does not contain, at a real cost in conditioning.

\textbf{Curvature is not identified.} The integrated yield carries no $\kappa^2$ term, so $\theta_{\kappa}$ recovers the autonomous yield and $\lambda_2$ is unidentified. Strict concavity is assumed, not tested. Any claim about the frontier's shape would require a $\kappa^2$ regressor this equation does not carry.

\textbf{The interaction span must be complete.} This requirement comes from the econometrics, not the theory, and violating it silently invalidated the earlier estimates. Centering a base series before forming a break-interacted cross-product expands as

\begin{equation}
(\kappa - \kappa_0)(\omega - \omega_0)D_t = \kappa\omega D_t - \kappa_0(\omega D_t) - \omega_0(\kappa D_t) + \kappa_0\omega_0 D_t
\end{equation}

so the design must span $\omega D_t$ and $\kappa D_t$ as well as $\kappa\omega D_t$ and $D_t$. Otherwise centering stops being a reparameterization and becomes a change of model. On synthetic data with a known data-generating process, ordinary least squares gives a relative spread across centering regimes of $6\times10^{-16}$ when the span is complete and $7.8\times10^{-2}$ when it is not. Two alternative explanations were tested and rejected: the estimator is not at fault (IM-OLS varies by about 0.3 percent on well-conditioned data) and neither is ill-conditioning (forcing the condition number to 181 leaves the spread at 0.26 percent). B5 carries the complete span.

\FloatBarrier
\section{Results}

The estimation window is 1929--2024 ($T = 94$). The internal break date $T_B^{*} = 1969$ comes from a Bai-Perron search on $\Delta(\kappa_t\omega_{t-1})$ and is not imposed.

\paragraph{What the integration-order requirement rules out.} A cointegrating regression is only interpretable as a long-run relationship if its variables share a common stochastic trend. The practical question is counterfactual: which specifications would this requirement \textit{reject}? Two kinds. A variable that is stationary around its breaks --- trend-stationary rather than integrated --- cannot participate in a cointegrating vector at all, because it has no stochastic trend to share; the lagged wage share fails this way in the aggregate specification of the companion report. A variable that remains integrated after first differencing carries a second unit root, and regressing it on $I(1)$ variables produces residuals that look stationary in-sample without any long-run relationship existing; the mechanization gap fails this way in the Post-Fordist sub-sample. Here, on the full sample, all three base variables survive both tests: each has a unit root in levels and loses it in first differences, tested under broken deterministics. The route is therefore \textbf{formal} rather than exploratory, and the estimates below are not carrying an unstated caveat about what the series are.

Step controls $D_{1933}$, $B_{1942\text{--}1945}$, $D_{1981}$ and $D_{2000}$ enter additively as variance hygiene. A reader concerned that these dummies are absorbing the regime shift rather than merely cleaning variance should note that the break-interaction decision below comes back negative: the additive controls carry what shift there is, and the multiplicative term adds nothing significant.

I report all three centering regimes and designate no single primary model. The reason is in Section~\ref{sec:deviations}: no configuration clears the condition-number threshold, and the paper-facing claim is the stability of the interaction coefficients across regimes, not any one point estimate.

\begin{table}[H]
\centering
\caption{Specification B5, Uninteracted ($T = 94$, 10 design columns, 83 residual d.f.)}
\label{tab:b5_nointer}
\small
\begin{tabular}{p{5.2cm}p{2.9cm}p{2.9cm}p{2.9cm}}
\toprule
Parameter & Uncentered & Initial-state & Sample-mean \\
\midrule
$\hat{\theta}_{NRC}$ & 0.326$^{***}$ (4.82) & 0.321$^{***}$ (4.48) & 0.371$^{***}$ (6.11) \\
$\hat{\theta}_{\kappa}$ & -3.130 (-1.00) & 0.980$^{***}$ (2.65) & 1.117$^{***}$ (4.74) \\
$\hat{\beta}_{p^{\kappa}}$ & 0.313 (0.56) & -0.807$^{**}$ (-2.35) & -0.266 (-1.22) \\
$\hat{\beta}_{\omega}$ & 1.606 (0.44) & -4.191$^{**}$ (-2.18) & -2.605$^{**}$ (-1.99) \\
$\hat{\theta}_{\omega\kappa}$ & 5.693 (1.17) & 4.203 (0.98) & 4.010 (0.87) \\
$\hat{\theta}_{\kappa p^{\kappa}}$ & 1.027 (1.44) & 0.997 (1.40) & 0.655 (0.89) \\
\midrule
Condition number & 638.06 & 49.50 & \textbf{30.02} \\
Shin LM ($C_{\text{Shin}}$) & 0.0523 & 0.0540 & 0.0610 \\
Residual ADF & -4.458$^{**}$ & -4.632$^{**}$ & -4.408$^{*}$ \\
$\hat{\sigma}_u$ & 0.1095 & 0.1094 & 0.1125 \\
$R^2_{\text{level}}$ & 0.9912 & 0.9912 & 0.9907 \\
\bottomrule
\multicolumn{4}{p{15.5cm}}{\footnotesize Note: Fixed-$b$ HAR $t$-statistics in parentheses. $^{*}p<0.10, ^{**}p<0.05, ^{***}p<0.01$. Source: \texttt{specB5\_coefficients.csv}.}
\end{tabular}
\end{table}

\begin{table}[H]
\centering
\caption{Specification B5, Break-Interacted at $T_B^{*} = 1969$ ($T = 94$, 14 design columns, 79 residual d.f.)}
\label{tab:b5_inter}
\small
\begin{tabular}{p{5.2cm}p{2.9cm}p{2.9cm}p{2.9cm}}
\toprule
Parameter & Uncentered & Initial-state & Sample-mean \\
\midrule
$\hat{\theta}_{NRC}$ & 0.593$^{***}$ (8.00) & 0.588$^{***}$ (8.09) & 0.599$^{***}$ (8.21) \\
$\hat{\theta}_{\kappa}$ & 7.430$^{*}$ (1.91) & 0.994$^{***}$ (3.16) & 0.944$^{***}$ (4.48) \\
$\hat{\beta}_{p^{\kappa}}$ & -0.430 (-0.98) & 0.221 (0.86) & -0.053 (-0.29) \\
$\hat{\beta}_{\omega}$ & -12.205$^{**}$ (-2.00) & -1.573 (-1.23) & -7.674$^{***}$ (-3.36) \\
$\hat{\theta}_{\omega\kappa}$ & -9.613 (-1.60) & -9.694 (-1.62) & -13.822$^{**}$ (-2.40) \\
$\hat{\theta}_{\kappa p^{\kappa}}$ & -0.617 (-1.15) & -0.558 (-1.01) & -0.689 (-1.24) \\
$\hat{\gamma}_1$ (level shift) & -11.195$^{**}$ (-2.22) & 0.677$^{**}$ (2.39) & 0.018 (0.19) \\
$\hat{\theta}_{\kappa,\text{brk}}$ & -10.417 (-1.56) & -1.235$^{***}$ (-2.87) & -1.384$^{***}$ (-3.49) \\
$\hat{\beta}_{\omega,\text{brk}}$ & 16.469$^{**}$ (2.05) & 0.817 (0.15) & 9.220$^{***}$ (3.12) \\
$\hat{\theta}_{\omega\kappa,\text{brk}}$ & 14.418 (1.34) & 14.139 (1.31) & 18.427$^{*}$ (1.75) \\
\midrule
Condition number & 2531.72 & 82.72 & 39.94 \\
Shin LM ($C_{\text{Shin}}$) & 0.0152 & 0.0146 & 0.0165 \\
Residual ADF & -7.400$^{***}$ & -7.366$^{***}$ & -7.053$^{***}$ \\
$\hat{\sigma}_u$ & 0.0641 & 0.0640 & 0.0658 \\
$R^2_{\text{level}}$ & 0.9971 & 0.9971 & 0.9970 \\
\bottomrule
\multicolumn{4}{p{15.5cm}}{\footnotesize Note: Fixed-$b$ HAR $t$-statistics in parentheses. $^{*}p<0.10, ^{**}p<0.05, ^{***}p<0.01$. Step controls $D_{1933}$, $B_{1942\text{--}1945}$, $D_{1981}$, $D_{2000}$ enter all columns. Source: \texttt{specB5\_coefficients.csv}.}
\end{tabular}
\end{table}

Cointegration holds throughout. $C_{\text{Shin}}$ runs from 0.0146 to 0.0610 against a 10 percent critical value of 0.095, and the residual ADF rejects no-cointegration in all six columns. The break interaction improves fit substantially, cutting $\hat{\sigma}_u$ from 0.109 to 0.064.

Centering changes two things by construction --- the conditioning of the design and the point at which baseline linear coefficients are evaluated. The invariance claim in the centering protocol therefore applies to the interaction-class parameters, and only those are gated.

\begin{table}[H]
\centering
\caption{Sensitivity of Each Parameter to the Centering Regime (Break-Interacted Model)}
\label{tab:centering}
\small
\begin{tabular}{p{4.2cm}p{2.2cm}p{4.0cm}p{2.0cm}p{1.7cm}}
\toprule
Parameter & Class & Range across regimes & Rel. spread & Sign stable \\
\midrule
$\theta_{NRC}$ & interaction & [0.588, 0.599] & 0.018 & Yes \\
$\theta_{\kappa p^{\kappa}}$ & interaction & [-0.689, -0.558] & 0.210 & Yes \\
$\theta_{\omega\kappa,\text{brk}}$ & interaction & [14.139, 18.427] & 0.274 & Yes \\
$\theta_{\omega\kappa}$ & interaction & [-13.822, -9.613] & 0.381 & Yes \\
\midrule
$\beta_{\omega}$ & base & [-12.205, -1.573] & 1.487 & Yes \\
$\theta_{\kappa,\text{brk}}$ & base & [-10.417, -1.235] & 2.113 & Yes \\
$\theta_{\kappa}$ & base & [0.944, 7.430] & 2.077 & Yes \\
$\beta_{p^{\kappa}}$ & base & [-0.430, 0.221] & 2.775 & No \\
$\gamma_1$ & base & [-11.195, 0.677] & 2.996 & No \\
\bottomrule
\multicolumn{5}{p{15.5cm}}{\footnotesize Note: Seven of seven interaction-class parameters hold their sign across all three centering regimes. Base-class parameters shift with the point of evaluation by construction and are not gated. Source: \texttt{specB5\_centering\_stability.csv}.}
\end{tabular}
\end{table}

\FloatBarrier
\section{Specification Space, and What Failed}

Reporting only the surviving specification would misrepresent how much of the space was searched. Four designs were estimated and superseded.

\textbf{B3, the canonical fully interacted split-sample model.} Carries $\kappa_t$ three times --- as a level, in $\kappa\omega$, and in $\kappa p^{\kappa}$ --- with condition numbers of 82 and 361 in the Fordist window and 43 and 231 in the Post-Fordist. Its break-interacted coefficients are centering-dependent: $\hat{\theta}_{\omega\kappa}$ spans $[-23.61, 47.79]$ across regimes in the interacted Fordist column.

\textbf{B4, the standalone-gap exclusion.} Dropped $\theta_{\kappa}$ on the argument that mechanization is distribution-conditioned rather than autonomous. Conditioning improved substantially (Fordist uninteracted CN 82.2 to 21.2, the one cell in the project that clears 30). But Section~4 shows the exclusion is not licensed by the derivation, and residual variance rose in every cell. Its $\hat{\theta}_{\omega\kappa}$ spans $[-19.60, 5.05]$.

\textbf{The span-broken full-sample track.} Correct in restoring the full 94-observation window and in locating the break endogenously, but it carried $\kappa\omega D_t$ without $\kappa D_t$ or $\omega D_t$. Its $\hat{\theta}_{\kappa p^{\kappa}}$ spans $[-0.744, 2.083]$ --- the sign of the price channel was set by the centering convention.

\textbf{Deterministic specification.} A linear trend added to the deterministic component was tested as a candidate remedy for the centering sensitivity and rejected: the median relative spread moved only from 1.98 to 1.64.

\begin{table}[H]
\centering
\caption{Sign Identification Across Tracks}
\label{tab:tracks}
\small
\begin{tabular}{p{5.6cm}p{3.2cm}p{1.6cm}p{3.0cm}p{1.6cm}}
\toprule
Track & $\hat{\theta}_{\omega\kappa}$ range & Stable & $\hat{\theta}_{\kappa p^{\kappa}}$ range & Stable \\
\midrule
Split-sample, Fordist interacted & [-23.607, 47.794] & No & [-11.482, 0.012] & No \\
Split-sample, Post-Fordist interacted & [-50.119, -45.959] & Yes & [0.463, 0.604] & Yes \\
Full-sample, span broken & [-19.600, 5.053] & No & [-0.744, 2.083] & No \\
\textbf{Full-sample, span complete (B5)} & \textbf{[-13.822, -9.613]} & \textbf{Yes} & \textbf{[-0.689, -0.558]} & \textbf{Yes} \\
\bottomrule
\multicolumn{5}{p{15.5cm}}{\footnotesize Note: Ranges are across the three centering regimes. Source: \texttt{specB5\_vs\_prior\_tracks\_comparison.csv}.}
\end{tabular}
\end{table}

\FloatBarrier
\section{Key Findings}

\textbf{The theory's two sign restrictions are not jointly satisfied.} Section~4 requires $\gamma_1 > 0$ and $\gamma_2 > 0$. Recovering them from the structural identities:

\begin{table}[H]
\centering
\caption{Structural Parameter Recovery Against the Theoretical Sign Restrictions}
\label{tab:signs}
\small
\begin{tabular}{p{3.9cm}p{3.2cm}p{2.2cm}p{3.2cm}p{2.2cm}}
\toprule
Specification & $\hat{\gamma}_1 = \hat{\theta}_{\omega\kappa}$ & Sign correct & $\hat{\gamma}_2 = -\hat{\theta}_{\kappa p^{\kappa}}$ & Sign correct \\
\midrule
Uninteracted & 4.203 & \textbf{Yes} & -0.997 & No \\
Break-interacted & -9.694 & No & 0.558 & \textbf{Yes} \\
\bottomrule
\multicolumn{5}{p{15.5cm}}{\footnotesize Note: Initial-state centering. Theory requires $\gamma_1 > 0$ (wage-savings dividend) and $\gamma_2 > 0$ (relative price penalty) jointly. Source: \texttt{specB5\_structural\_parameter\_recovery.csv}.}
\end{tabular}
\end{table}

Each specification recovers exactly one channel of the viability condition and reverses the other. Read as regime signatures, this is not a technical failure but a substantive one: the induced-innovation mechanism that Section~\ref{sec:theory} derives does not appear in the US data as a single coherent structure over 1929--2024. The descriptive evidence in Section~2 points the same way, with $\kappa$ and $p^{\kappa}$ co-moving positively in both eras where the theory requires them to move oppositely.

\textbf{The price channel is now sign-identified, but not significant.} $\hat{\theta}_{\kappa p^{\kappa}}$ holds a negative sign across all three centering regimes, $[-0.689, -0.558]$, where every earlier track flipped with the convention chosen. Its $t$-statistic is $-1.01$. The price penalty is directionally consistent with Section 3.2 and statistically indistinguishable from zero.

\textbf{Extensive scale is the one sharply identified parameter, and it carries an overaccumulation punchline.} $\hat{\theta}_{NRC} = 0.588$ ($t = 8.09$), with a relative spread of 0.018 across regimes. Plant structures anchor capacity precisely, and the failure documented above is specific to the state-dependent channel.

Put the elasticity together with the descriptive rates in Table~\ref{tab:descriptive} and the stake becomes concrete. After 1974 structures accumulated at 5.72 percent a year. At an elasticity of 0.588, the extensive margin alone contributed

\begin{equation}
0.588 \times 5.72 = 3.36 \ \text{percent per year}
\end{equation}

of capacity growth, against realized output growth of 3.12 percent. Capacity from plant structures alone outran output by roughly 0.24 points a year, which compounds to about 12 percentage points of accumulated slack across the half-century --- and that is before any contribution from machinery. This is the arithmetic behind the overaccumulation reading: it does not require the state-dependent channel to work, only the extensive elasticity to be below one and capital to keep accumulating.

\textbf{The elasticity is probably biased downward.} $k^{NRC}_t$ is a perpetual-inventory construct depending on assumed service lives and depreciation geometry, so it carries classical measurement error on the right-hand side. In a cointegrating regression the attenuation is less severe than in a stationary one, since the super-consistency of the long-run estimator dominates as $T$ grows, but it does not vanish: the direction is toward zero, so 0.588 should be read as a lower bound on the true extensive elasticity. That direction matters for the punchline above. A larger true $\theta_{NRC}$ widens the gap between capacity growth and output growth, strengthening rather than weakening the overaccumulation reading.

\textbf{The 1969 break does not survive the span correction.} $\hat{\theta}_{\omega\kappa,\text{brk}} = 14.139$ with $t = 1.31$, so the decision is to retain the additive step control: no significant multiplicative shift in the wage-mechanization relation. An earlier specification reported this coefficient at $t = -2.90$, and that significance was an artifact of the incomplete span.

\textbf{Two parameters are unrecoverable.} $\lambda_2$ requires a $\kappa^2$ regressor the equation does not carry, so strict concavity is assumed rather than tested. The transformation elasticity $\theta^{SpecB}_t$ is not re-derived here; it runs through the LOESS path-derivative stage rather than the break-aware one, and re-deriving it under B5 requires re-estimating the full model grid. No B5-conditioned trajectory is asserted.

\FloatBarrier
\section{Deviations from Standard Protocol}
\label{sec:deviations}

Three departures from the applied-econometrics protocol used elsewhere in this project, each with its reason.

\textbf{Frisch-Waugh-Lovell orthogonalization of interaction terms is not performed.} The general protocol calls for substituting the residual of an auxiliary regression for the raw cross-product when $I(2)$ risk is present. The E10V estimator lock forbids this, and the argument is specific to the estimator: partial summation does not commute with projection ($S M_X \neq M_{SX} S$), so level-filtered residuals violate the Brownian limit theory that cointegrating polynomial regression relies on. FWL pre-filtering is not useful in this time-series setting. The interaction enters raw, and \texttt{s33\_assert\_no\_fwl()} enforces it at runtime.

\textbf{Integration order is established by Lee-Strazicich rather than ADF and KPSS.} The general protocol screens with ADF and KPSS on levels and first differences. Every series here contains structural breaks, and the augmented Dickey-Fuller test is biased toward non-rejection under a break --- it would systematically misclassify exactly the series this chapter is about. The Lee-Strazicich Minimum-LM test with endogenous break location replaces it, adjudicated on the levels-and-differences pair. Cointegration diagnostics retain the residual ADF and the Shin LM statistic.

\textbf{The condition-number threshold is not met, and no model is designated primary.} The protocol blocks any specification with a standardized condition number above 30 regardless of theoretical motivation. No B5 configuration clears it: the uninteracted model reaches 30.02 under sample-mean centering, and the break-interacted model runs from 39.94 to 2531.72. Rather than select the regime that comes closest and present it as preferred, this report gives the full triad and rests its claims on sign stability across regimes. B5 is entered below as \textbf{Held}, not Promoted.

\textbf{Bootstrap inference is not triggered.} Small-sample bootstrap applies at $N \le 30$; here $T = 94$, so asymptotic Fixed-$b$ HAR inference stands alone.

\FloatBarrier
\begin{landscape}
\section{Consolidated Verdict Table}

\begin{table}[H]
\centering
\caption{Consolidated Specification Verdict}
\label{tab:verdict}
\small
\begin{tabular}{p{3.6cm}p{4.4cm}p{6.4cm}p{2.2cm}p{6.0cm}}
\toprule
Candidate Variable & Theoretical Basis & Integration \& Collinearity Evidence & Status & Governing Lock / Rationale \\
\midrule
$k^{NRC}_t$ & Extensive capacity envelope, \S3.2.1 & LS $I(1)$-with-break, gate PASS; $\hat{\theta}_{NRC}$ rel.\ spread 0.018 & \textbf{Promoted} & Clears integration gate; sharply identified across all regimes \\
$\kappa_t$ & Autonomous yield $(\lambda_1-\gamma_0)$, \S3.2.3 & LS $I(1)$-with-break, gate PASS; base-class, rel.\ spread 2.077 & \textbf{Held} & Required by the derivation; centering-sensitive, so reported as a range \\
$p^{\kappa}_t$ & Relative price penalty $\gamma_2$, \S3.2.2 & LS $I(1)$-with-break, gate PASS; $\hat{\beta}_{p^{\kappa}}$ sign-unstable & \textbf{Held} & Enters as a level control; the identified channel is the interaction \\
$\omega_{t-1}$ & Wage-savings dividend $\gamma_1$, \S3.2.2 & Level term sign-stable but rel.\ spread 1.487 & \textbf{Held} & Base-class; point of evaluation moves with centering by construction \\
$\kappa_t \cdot \omega_{t-1}$ & $\theta_{\omega\kappa}=\gamma_1>0$, \S3.2.3 & Sign-stable $[-13.82, -9.61]$; \textbf{sign contradicts theory} & \textbf{Held} & Identified but sign-reversed in the interacted model; H1 not supported \\
$\kappa_t \cdot p^{\kappa}_t$ & $\theta_{\kappa p\kappa}=-\gamma_2<0$, \S3.2.3 & Sign-stable $[-0.689, -0.558]$; $t=-1.01$ & \textbf{Held} & Sign consistent with \S3.2, magnitude indistinguishable from zero \\
$\kappa_t\omega_{t-1}D_t$ (1969) & Regime shift in the wage channel, \S2.6 & $t=1.31$; sign-stable $[14.14, 18.43]$ & \textbf{Held} & Not significant at 5\%; decision is additive step control \\
$\kappa_t\omega_{t-1}p^{\kappa}_t$ & None --- $MC$ additively separable & Not estimated & \textbf{Rejected} & \S3.2 derives no three-way term; would add collinearity without warrant \\
$\kappa^2_t$ & Frontier curvature $\lambda_2$ & Not in the integrated $\Psi_t$ & \textbf{Blocked} & Unidentified; strict concavity assumed, not tested \\
FWL-residualized interactions & Protocol Gate 4 & Not estimated & \textbf{Blocked} & E10V: partial summation does not commute with projection \\
\midrule
\textbf{Specification B5 (whole)} & \S3.2 Eq.~\ref{eq:b5} & Cointegrated ($C_{\text{Shin}}$ 0.015--0.061); \textbf{CN 30.0--2531.7 fails the $\le$30 gate} & \textbf{Held} & Sign-identified across centering, but no configuration clears the collinearity threshold \\
\bottomrule
\end{tabular}
\end{table}
\end{landscape}

\FloatBarrier
\section{Open Items}

\begin{itemize}
\item The \S3.2 versus A06 sign conflict on the price channel is unresolved and now bears directly on interpretation, since the descriptive co-movement sits with A06 while the estimate sits with \S3.2.
\item The split-sample and span-broken tracks retain incomplete interaction spans; their break-interacted coefficients should not be quoted for either hypothesis until corrected. The v7 report's Table 3.2 $\delta$-rows and Table 3.3 columns rest on those designs.
\item Residual interaction-coefficient spread of 0.21 to 0.42 exceeds the roughly 0.007 that a complete span achieves on synthetic data. B5 is sign-identified, not exactly invariant, and the remaining gap is unexplained.
\item $\theta^{SpecB}_t$ awaits re-derivation under B5 through the path-derivative stage.
\item No configuration clears the condition-number threshold. A specification that does --- plausibly one built on investment flows rather than stock ratios --- remains the outstanding econometric problem.
\end{itemize}

\end{document}
